\section{How each literature tests}
\label{sec:testing}

A test across the two literatures puts one side's fix in front of the
other's problem. Whether that test is fair depends on how each fix was built
and evaluated, so we coded this for all \NDesignFixes{} fixes in the corpus
that act on retrieved text. For each we recorded, from the full text, how
many passages one test item hands the fix, what the misleading text competes
with, whether results are split by what the model already knew, and whether
the attacker adapts to the fix. Every coded field carries a verbatim quote,
\PctDesignQuotes{}\% of which an automated check finds in the paper's text,
and a hand check of \NDesignHand{} papers agreed on the passage count in
every case. \NDesignInScope{} of the fixes test with misleading text,
\NDesignSec{} from security papers and \NDesignSaf{} from safety papers; the
counts below are over those.

\begin{table}[t]
\centering
\caption{What each fix is tested on, by the question it asks, over the
\NDesignInScope{} retrieval fixes whose evaluation contains misleading text.}
\label{tab:units}
\footnotesize
\begin{tabular}{@{}p{0.34\columnwidth}p{0.24\columnwidth}p{0.32\columnwidth}@{}}
\toprule
\textbf{Question the fix asks} & \textbf{Tested on} & \textbf{Fixes} \\
\midrule
Document or memory? & one passage & \NDesignMemOne{} of the
  \NDesignMemCounted{} safety fixes that state a count \\
Which passage is right? & a set of passages & \NDesignSetSec{} security,
  \NDesignSetSaf{} safety \\
Is there an injected instruction? & one text & \NDesignInstrOne{} of the
  \NDesignInstrCounted{} security fixes that state a count \\
\bottomrule
\end{tabular}
\end{table}

\paragraph{The number of passages follows the question, not the literature.}
Table~\ref{tab:units} groups the fixes by what they decide. A fix that asks
whether to trust the document or the model's own memory is tested on one
passage, because one is enough to set against memory; every such fix comes
from a safety paper. A fix that asks which retrieved passage to trust is
tested on a set, in both literatures: the security defenses that vote across
passages and the safety fixes for noisy or conflicting retrieval. Detectors
for injected instructions read one text at a time. Where a set is used, its
median size is \NDesignMedianK{} on both sides.

\paragraph{Three choices separate the literatures.} What does split them is
what the bad text is assumed to do. No security fix sets the misleading text
against the model's memory (\NDesignMemSec{} of \NDesignSec{}); in the
safety literature this is the most common design. Only safety papers report
results split by whether the model already knew the answer
(\NDesignSplitSaf{} against \NDesignSplitSec{}). Only security papers test an
attacker who knows the fix (\NDesignAdaptSec{} against
\NDesignAdaptSaf{}). A document that has gone out of date, the standard example of an honest
mistake, is almost never tested with real stale text: \NDesignOutdated{} of the \NDesignSaf{} safety fixes do, and no
security fix does.

\paragraph{What a test across the literatures needs.} A fix has to receive
the input its question requires, or the test measures the format and not the
method. Our first cross tests did not meet this. They joined ten retrieved
passages into one long document for the safety fixes built on one passage,
and some fixes broke on that input: the released code of the safety fix
COIECD cuts prompts beyond 2,048 tokens from the end, which removes the
question, and our runner shortened the safety fix SpARE's context so far that
the planted passages were cut. We withdrew those results. What remains, and
what Section~\ref{sec:results} reports, runs every safety fix on one
passage, its own format. The reverse direction needs what neither literature
has: ranked sets that contain real honest mistakes, with a per-model record
of what the model already knew. Section~\ref{sec:dataset} builds it.

Answer to RQ2. Each fix is tested on the input its question needs: one
passage when it weighs a document against memory, a set when it chooses
among passages. The literatures differ in what the bad text competes with,
in whether they split results by prior knowledge, and in whether the
attacker adapts. A fair test across them gives each fix its own format and
reports both literatures' measures.
